% Luc Destombe --- licence CC BY-NC-SA 4.0
% https://creativecommons.org/licenses/by-nc-sa/4.0/deed.fr
% Fichier autonome : compiler avec pdflatex, deux fois (signets, nombre de pages).
\documentclass[a4paper,10pt]{article}
\makeatletter
\newif\iflycee@niveaux
\newif\iflycee@palatino
\lycee@palatinotrue

\RequirePackage[utf8]{inputenc}
\RequirePackage[T1]{fontenc}
\RequirePackage[french]{babel}
\RequirePackage{etoolbox}

\newcommand{\ShorthandsOff}{\@ifundefined{shorthandoff}{}{\shorthandoff{;:!?}}}
\newcommand{\ShorthandsOn}{\@ifundefined{shorthandon}{}{\shorthandon{;:!?}}}
\ShorthandsOff
\AtBeginDocument{\ShorthandsOn}
\AtBeginEnvironment{tikzpicture}{\ShorthandsOff}

\RequirePackage{geometry}
\RequirePackage{fancyhdr}
\RequirePackage{lastpage}
\RequirePackage{multicol}
\RequirePackage{array}
\RequirePackage{enumitem}
\RequirePackage{xcolor}
\RequirePackage{needspace}

\RequirePackage{amsmath,amssymb}
\RequirePackage{mathpazo}
\DeclareMathSizes{10.95}{10.95}{8}{5.5}
\begingroup\catcode`\_=\active \catcode`\^=\active
\gdef_#1{\sb{\scriptscriptstyle #1}}
\gdef^#1{\sp{\scriptscriptstyle\lycee@moinsepais #1}}
\endgroup
\newcommand\lycee@moins{\mathbin{%
  \setbox\z@\hbox{$\m@th\scriptscriptstyle\lycee@moinsorig$}%
  \dimen@=.55\wd\z@ \advance\dimen@-.5pt
  \hbox to.75\wd\z@{\hss$\m@th\scriptscriptstyle\vcenter{\hbox{%
    \tikz\draw[line width=.5pt,line cap=round](0,0)--(\the\dimen@,0);}}$\hss}}}
\begingroup\lccode`\~=`\- \lowercase{\endgroup\let~}\lycee@moins
\newcommand\lycee@moinsepais{\mathcode`\-="8000 }
\AtBeginDocument{\mathchardef\lycee@moinsorig=\mathcode`\-\relax}
\AtBeginDocument{\catcode`\_=12 \mathcode`\_="8000
  \catcode`\^=12 \mathcode`\^="8000 }
\newcommand\lycee@exposants{%
  \fontdimen13\textfont2=.48\fontdimen6\textfont2
  \fontdimen14\textfont2=.48\fontdimen6\textfont2
  \fontdimen15\textfont2=.48\fontdimen6\textfont2
  \fontdimen13\scriptfont2=.48\fontdimen6\scriptfont2
  \fontdimen14\scriptfont2=.48\fontdimen6\scriptfont2
  \fontdimen15\scriptfont2=.48\fontdimen6\scriptfont2}
\every@math@size\expandafter{\the\every@math@size\lycee@exposants}
\newlength{\retraitcalculs}
\setlength{\retraitcalculs}{2em}

\RequirePackage{tikz}
\usetikzlibrary{arrows.meta,calc,babel,shapes.misc}
\RequirePackage[most]{tcolorbox}
\tcbuselibrary{skins,breakable}

\definecolor{coulDef}{HTML}{1F4E79}
\definecolor{coulProp}{HTML}{7030A0}
\definecolor{coulRem}{HTML}{C55A11}
\definecolor{coulEx}{HTML}{2E7D32}
\definecolor{coulExo}{HTML}{B03A2E}
\definecolor{coulPart}{HTML}{1F4E79}
\definecolor{coulMeth}{HTML}{1B6E4F}
\definecolor{coulCourbe}{HTML}{0B5ED7}
\definecolor{coulCourbeB}{HTML}{C00000}
\definecolor{coulCourbeC}{HTML}{2E7D32}
\definecolor{coulGrille}{HTML}{8C8C8C}
\definecolor{coulRep}{HTML}{1B6E4F}
\definecolor{coulBar}{HTML}{2E7D32}

\newcommand{\R}{\mathbb{R}}
\newcommand{\Rs}{\mathbb{R}^{*}}
\renewcommand{\le}{\leqslant}
\renewcommand{\ge}{\geqslant}
\newcommand{\vect}[1]{\overrightarrow{#1}}
\newcommand{\norme}[1]{\left\lVert #1 \right\rVert}
\newcommand{\intervalle}[2]{\left[#1\,;#2\right]}
\RequirePackage{cancel}
\renewcommand{\CancelColor}{\color{coulExo}}
\newcommand{\fois}[1]{\times{\color{coulExo}#1}}
\newcommand{\barre}[1]{\cancel{#1}}

\tikzset{grille/.style={coulGrille,line width=0.4pt}}
\newcommand{\quadrillage}[4]{\draw[grille,step=1] (#1,#2) grid (#3,#4);}
\tikzset{
  croix/.style={cross out,draw,line width=0.6pt,inner sep=0pt,minimum size=4pt},
  croixdroite/.style={croix,rotate=45,line width=0.8pt},
}
\newcommand{\pt}[3]{\path (#1) node[croix]{} node[#2,font=\small,inner sep=2.5pt]{$#3$};}

\newcommand{\lycee@repere}[4]{%
  \draw[grille,step=1] (#1,#2) grid (#3,#4);
  \draw[-{Stealth[length=2mm]},thick] (#1,0)--({#3+0.4},0) node[below left=-1pt]{$x$};
  \draw[-{Stealth[length=2mm]},thick] (0,#2)--(0,{#4+0.4}) node[below left=-1pt]{$y$};
  \node[below left,font=\scriptsize,inner sep=1.5pt] at (0,0) {$0$};}
\newcommand{\lycee@gradx}[1]{\foreach \k/\l in {#1}{%
  \draw (\k,0.07)--(\k,-0.07) node[below,font=\scriptsize,inner sep=1.5pt]{$\l$};}}
\newcommand{\lycee@grady}[1]{\foreach \k/\l in {#1}{%
  \draw (0.07,\k)--(-0.07,\k) node[left,font=\scriptsize,inner sep=1.5pt]{$\l$};}}
\newcommand{\lycee@intff}[2]{\left[#1\,;#2\right]}
\newcommand{\lycee@intfo}[2]{\left[#1\,;#2\right[}
\newcommand{\lycee@intof}[2]{\left]#1\,;#2\right]}
\newcommand{\lycee@intoo}[2]{\left]#1\,;#2\right[}
\newcommand{\lycee@ens}[1]{\left\{#1\right\}}
\AtBeginDocument{%
  \providecommand{\repere}{\lycee@repere}%
  \providecommand{\gradx}{\lycee@gradx}%
  \providecommand{\grady}{\lycee@grady}%
  \providecommand{\Cf}{\mathcal{C}\sb{\scriptscriptstyle f}}%
  \providecommand{\Cg}{\mathcal{C}\sb{\scriptscriptstyle g}}%
  \providecommand{\intff}{\lycee@intff}%
  \providecommand{\intfo}{\lycee@intfo}%
  \providecommand{\intof}{\lycee@intof}%
  \providecommand{\intoo}{\lycee@intoo}%
  \providecommand{\ens}{\lycee@ens}}

\newlist{questions}{enumerate}{3}
\setlist[questions,1]{label=\arabic*\textdegree),leftmargin=*,itemsep=2pt,topsep=3pt}
\setlist[questions,2]{label=\alph*),leftmargin=*,itemsep=1pt,topsep=2pt}
\setlist[questions,3]{label=\roman*),leftmargin=*,itemsep=1pt,topsep=1pt}
\setlist[itemize]{label=\textbullet,leftmargin=*,itemsep=2pt,topsep=3pt}

\newcommand{\besoinplace}[1]{\par\penalty\@M\begingroup
  \dimen@=#1\relax\dimen@ii=\pagegoal\advance\dimen@ii-\pagetotal
  \ifdim\dimen@>\dimen@ii\ifdim\dimen@ii>\z@\vfil\fi\break\fi\endgroup}

\newcommand{\lycee@pied}{\thepage/\pageref{LastPage}}
\newcommand{\entete}[2]{%
  \pagestyle{fancy}\fancyhf{}%
  \fancyhead[L]{\footnotesize\color{coulPart}\textbf{#1}}%
  \fancyhead[R]{\footnotesize\color{coulPart}\textbf{#2}}%
  \fancyfoot[C]{\footnotesize\lycee@pied}%
  \renewcommand{\headrulewidth}{0.4pt}%
  \renewcommand{\headrule}{\hbox to\headwidth{%
    \color{coulPart!40}\leaders\hrule height \headrulewidth\hfill}}}

\RequirePackage{ccicons}
\newcommand{\lycee@licence}{{\scriptsize\color{gray}%
  \copyright\ Luc Destombe --- \ccbyncsa\ CC BY-NC-SA 4.0}}
\newif\iflycee@licence \lycee@licencetrue
\newcommand{\sanslicence}{\lycee@licencefalse}
\AtBeginDocument{\iflycee@licence\AddToHookNext{shipout/foreground}{%
  \put(0,0){\rlap{\kern\dimexpr1in+\hoffset+\oddsidemargin+\textwidth\relax
    \llap{\raisebox{-\dimexpr1in+\voffset+\topmargin+\headheight+\headsep
      +\textheight+\footskip\relax}{\lycee@licence}}}}}\fi}

\newcommand{\formulesserrees}{%
  \setlength{\abovedisplayskip}{3pt}\setlength{\belowdisplayskip}{3pt}%
  \setlength{\abovedisplayshortskip}{2pt}\setlength{\belowdisplayshortskip}{3pt}}

\setlength{\parindent}{0pt}

\RequirePackage[implicit=false,hidelinks,pdfencoding=auto,psdextra,bookmarksopen,
  bookmarksopenlevel=1,pdfcreator={},pdfproducer={}]{hyperref}
\RequirePackage{bookmark}
\hypersetup{pdfauthor={Luc Destombe},pdfkeywords={CC BY-NC-SA 4.0}}
\newcounter{lycee@signet}
\newcommand{\signet}[3][]{\stepcounter{lycee@signet}%
  \edef\lycee@dest{\if\relax\detokenize{#1}\relax
    signet.\arabic{lycee@signet}\else#1\fi}%
  \hypertarget{\lycee@dest}{}\bookmark[level=#2,dest=\lycee@dest]{#3}}
\pdfstringdefDisableCommands{%
  \def\R{R}\def\vect#1{#1}\def\Cf{Cf}\def\Cg{Cg}\def\fois#1{×#1}%
  \def\barre#1{#1}\def\intervalle#1#2{[#1;#2]}\def\intff#1#2{[#1;#2]}%
  \def\textsuperscript#1{#1}\def\dfrac#1#2{#1/#2}\def\frac#1#2{#1/#2}%
  \def\vec#1{vecteur #1}}
\RequirePackage[official]{eurosym}

\geometry{top=0.8cm,bottom=1.3cm,left=1cm,right=1cm,footskip=0.6cm}

\pagestyle{fancy}
\fancyhf{}
\renewcommand{\headrulewidth}{0pt}
\fancyfoot[C]{\footnotesize Page \thepage{} / \pageref{LastPage}}

\newcommand{\titrefiche}[2]{%
  \begin{center}
    \Large\textbf{#1}\\[2pt]
    \large\textbf{#2}\\[2pt]
  \end{center}
  \vspace{0.10cm}}

\setlength{\columnseprule}{0.4pt}
\setlength{\columnsep}{15pt}
\renewcommand{\columnseprulecolor}{\color{black!45}}
\setlength{\multicolsep}{2pt}

\newcounter{partiefiche}
\renewcommand{\thepartiefiche}{\Roman{partiefiche}}
\newif\ifapresbandeau
\newcommand{\lycee@bandeau}[1]{%
  \par\penalty-600
  \vspace{0.08cm}%
  \noindent\colorbox{coulPart}{%
    \parbox{\dimexpr\linewidth-2\fboxsep\relax}{%
      \color{white}\bfseries\raggedright\strut#1\strut}}%
  \nopagebreak\par\nopagebreak
  \vspace{0.06cm}\nopagebreak
  \apresbandeautrue}
\newcommand{\partiefiche}[1]{%
  \refstepcounter{partiefiche}%
  \lycee@bandeau{\signet{1}{\thepartiefiche{} --- #1}\thepartiefiche{} --- #1}}
\newcommand{\partiefichesansnum}[1]{\lycee@bandeau{\signet{1}{#1}#1}}

\newcounter{exo}
\newlength{\espaceexo}\setlength{\espaceexo}{7pt}
\newlength{\espacetitre}\setlength{\espacetitre}{2pt}
\newcommand{\exo}[1][\relax]{%
  \par
  \ifapresbandeau\apresbandeaufalse\else\penalty-150\addvspace{\espaceexo}\fi
  \refstepcounter{exo}%
  \noindent\signet[exercice-\arabic{exo}]{2}{Exercice \arabic{exo}}%
  \textbf{\color{coulExo}\underline{Exercice \theexo}}%
  \ifx\relax#1\relax\else\textbf{ : #1}\fi
  \par\nopagebreak\vspace{\espacetitre}\nopagebreak
  \ignorespaces}

\setlist[enumerate]{nosep,leftmargin=*,itemsep=2pt,topsep=2pt}
\setlist[itemize]{nosep,leftmargin=*,topsep=2pt}
\newcommand{\choix}[3]{\par\hspace*{1em}%
  \makebox[0.3\linewidth][l]{a) #1}\makebox[0.3\linewidth][l]{b) #2}c) #3}
\newcommand{\deux}[2]{\par\hspace*{0.6em}%
  \makebox[0.48\linewidth][l]{#1}#2}
\newcommand{\trois}[3]{\par\hspace*{0.6em}%
  \makebox[0.32\linewidth][l]{#1}\makebox[0.32\linewidth][l]{#2}#3}

\renewenvironment{center}{\par\addvspace{3pt}\centering}{\par\addvspace{3pt}}
\formulesserrees
\AtBeginDocument{\formulesserrees}
\clubpenalty=10000
\widowpenalty=10000
\displaywidowpenalty=10000
\brokenpenalty=10000
\setlength{\parskip}{0pt}

\RequirePackage{ragged2e}
\setlength{\RaggedRightRightskip}{0pt plus 1fil}
\AtBeginDocument{\RaggedRight\binoppenalty=10000 \relpenalty=10000 }
\g@addto@macro\@parboxrestore{\RaggedRight}
\makeatother

\newenvironment{expressions}[1][2]%
  {\par\addvspace{3pt}\renewcommand{\arraystretch}{#1}%
   \tabular{@{}p{0.47\linewidth}p{0.47\linewidth}@{}}}%
  {\endtabular\par\addvspace{3pt}}

\begin{document}

\titrefiche{Seconde --- Fractions et puissances}{Fiche d'exercices du chapitre}

\begin{multicols}{2}\raggedcolumns

\partiefiche{Nombres relatifs et priorités}

\exo[Calculs sur les relatifs]
Calculer :
\begin{expressions}
  $A=-9-7$ & $B=14-(-11)$\\
  $C=-6\times 7$ & $D=(-8)\times(-5)$\\
  $E=\dfrac{-(-27)}{-9}$ & $F=-13+(-4)$
\end{expressions}

\exo[Priorités de calcul]
Calculer en respectant les priorités :
\begin{expressions}
  $A=8^{2}-3\times(-6)$ & $B=-7-4^{2}$\\
  $C=6+(-9+5)\times 4$ & $D=-36\div(-4)\times 2^{2}$
\end{expressions}

\exo[Parenthèses]
Calculer en respectant les priorités :
\begin{expressions}
  $A=20-4\times(3-8)$ & $B=5\times(-3)+2^{3}$\\
  $C=-48\div\bigl(2\times(-3)\bigr)$ & $D=3-2\times\bigl(5-(1-6)\bigr)$
\end{expressions}

\partiefiche{Additionner et soustraire des fractions}

\exo[Même dénominateur]
Calculer et donner le résultat sous forme simplifiée :
\begin{expressions}
  $A=\dfrac{4}{5}+\dfrac{3}{5}$ & $B=\dfrac{-2}{9}+\dfrac{7}{9}$\\
  $C=\dfrac{5}{7}-\dfrac{9}{7}$ & $D=\dfrac{-3}{11}-\dfrac{-8}{11}$\\
  $E=\dfrac{7}{12}+\dfrac{-1}{12}$ & $F=\dfrac{-7}{4}-\dfrac{5}{4}$
\end{expressions}

\exo[Dénominateurs différents]
Calculer et donner le résultat sous forme simplifiée :
\begin{expressions}
  $A=\dfrac{2}{5}+\dfrac{3}{4}$ & $B=\dfrac{5}{6}+\dfrac{1}{4}$\\
  $C=\dfrac{3}{4}+\dfrac{1}{6}$ & $D=\dfrac{3}{8}-\dfrac{5}{6}$\\
  $E=-\dfrac{2}{3}-\dfrac{4}{5}$ & $F=-\dfrac{7}{6}+\dfrac{5}{4}$
\end{expressions}

\exo[Avec un nombre entier]
Calculer et donner le résultat sous forme simplifiée :
\begin{expressions}
  $A=3-\dfrac{5}{4}$ & $B=-2-\dfrac{3}{5}$\\
  $C=-4+\dfrac{17}{6}$ & $D=\dfrac{9}{2}-4$
\end{expressions}

\exo[Avec une lettre]
$x$ désigne un nombre. Écrire chaque expression avec une seule fraction :
\begin{expressions}
  $A=3x-\dfrac{x}{4}$ & $B=\dfrac{2}{5}x-3$\\
  $C=2-\dfrac{x-1}{3}$ & $D=\dfrac{2x+1}{3}-x$
\end{expressions}

\partiefiche{Multiplier et diviser des fractions}

\exo[Simplifier avant de calculer]
Simplifier, puis calculer, sans effectuer les produits :
\begin{expressions}
  $A=\dfrac{-4\times 9}{6\times 10}$ & $B=\dfrac{24\times(-35)}{(-21)\times 10}$\\
  $C=\dfrac{26\times(-9)}{15\times 39}$ & $D=\dfrac{14\times(-15)}{(-10)\times 21}$
\end{expressions}

\exo[Simplifier avec des lettres]
Simplifier ($a$ et $b$ sont des nombres non nuls) :
\begin{expressions}
  $A=\dfrac{15\times 8\times 14}{35\times 6\times 4}$ & $B=\dfrac{3x\times 20}{4\times 5}$\\
  $C=\dfrac{6a\times 5b}{10b\times 9}$ & $D=\dfrac{4a\times 9}{6\times 2a}$
\end{expressions}

\exo[Produits]
Calculer et donner le résultat sous forme simplifiée :
\begin{expressions}
  $A=\dfrac{5}{4}\times\dfrac{-3}{7}$ & $B=\dfrac{-2}{5}\times\dfrac{4}{3}$\\
  $C=-5\times\dfrac{3}{8}$ & $D=\dfrac{-7}{-3}\times\dfrac{2}{-5}$\\
  $E=\dfrac{-4}{5}\times\dfrac{-4}{9}$ & $F=\dfrac{3}{-8}\times\dfrac{-5}{-7}$
\end{expressions}

\exo[Quotients]
Calculer et donner le résultat sous forme simplifiée :
\begin{expressions}
  $A=\dfrac{2}{5}\div\dfrac{3}{7}$ & $B=\dfrac{5}{4}\div\dfrac{7}{3}$\\
  $C=\dfrac{3}{4}\div 6$ & $D=6\div\dfrac{-2}{5}$\\
  $E=\dfrac{4}{9}\div\dfrac{-5}{2}$ & $F=3\div\dfrac{-4}{7}$
\end{expressions}

\partiefiche{Fractions et priorités de calcul}

\exo[Respecter les priorités]
Calculer et donner le résultat sous forme simplifiée :
\begin{expressions}
  $A=\dfrac{5}{12}-\dfrac{-3}{4}$ & $B=\dfrac{7}{10}-\dfrac{-2}{5}$\\
  $C=-\dfrac{5}{6}+\dfrac{3}{4}-\dfrac{2}{3}$ & $D=\dfrac{9}{14}+\dfrac{5}{7}\times\dfrac{7}{10}$
\end{expressions}

\exo[Parenthèses et lettres]
Calculer et donner le résultat sous forme simplifiée :
\begin{expressions}[2.4]
  $A=\dfrac{4}{5}\left(\dfrac{1}{2}-\dfrac{1}{4}\right)$ &
  $B=\left(\dfrac{2}{3}-\dfrac{1}{6}\right)\times\left(\dfrac{1}{4}+\dfrac{1}{2}\right)$\\
  $C=\dfrac{3x}{4}-\dfrac{x-2}{6}$ & $D=\dfrac{3x}{4}\times\dfrac{8}{9}x$
\end{expressions}

\exo[Fractions de fractions]
Calculer et donner le résultat sous forme simplifiée :
\begin{expressions}[4.4]
  $A=\dfrac{\;\dfrac{3}{4}\;}{\;\dfrac{5}{6}\;}$ &
  $B=\dfrac{\;\dfrac{6}{-35}\;}{\;\dfrac{-9}{14}\;}$\\
  $C=\dfrac{\;\dfrac{7}{4}\;}{-14}$ &
  $D=\dfrac{\;2-\dfrac{3}{5}\;}{\;\dfrac{3}{4}-\dfrac{2}{5}\;}$
\end{expressions}

\partiefiche{Produit en croix}

\exo[Justifier le produit en croix]
Le nombre $x$ vérifie $\;\dfrac{3x-1}{4}=\dfrac{x+3}{2}$.
\begin{enumerate}[label=\arabic*)]
  \item Expliquer pourquoi $x$ vérifie aussi $2\times(3x-1)=4\times(x+3)$.
  \item En déduire la valeur de $x$.
\end{enumerate}

\exo[Premières équations]
Résoudre par le produit en croix :
\begin{expressions}[2.4]
  a) $\dfrac{x}{4}=\dfrac{5}{3}$ & b) $\dfrac{4-x}{3}=\dfrac{x+2}{5}$\\
  c) $\dfrac{2+x}{-3}=\dfrac{5}{6}$ & d) $\dfrac{2x-1}{3}=\dfrac{x+4}{-2}$
\end{expressions}

\exo[S'entraîner]
Résoudre, puis vérifier la solution :
\begin{expressions}[2.4]
  a) $\dfrac{x}{5}=\dfrac{3}{8}$ & b) $\dfrac{4}{9}=\dfrac{x}{6}$\\
  c) $\dfrac{x+2}{5}=\dfrac{2x-3}{4}$ & d) $\dfrac{2x+1}{3}=\dfrac{3x-2}{4}$
\end{expressions}

\exo[Plusieurs fractions (1)]
Résoudre en multipliant chaque terme par un dénominateur commun :
\begin{expressions}[2.4]
  a) $x+\dfrac{x}{5}=18$ & b) $\dfrac{x}{3}-\dfrac{x}{4}=5$\\
  c) $\dfrac{t}{3}-\dfrac{t+1}{5}=1$ & d) $\dfrac{3t-1}{2}+\dfrac{t-2}{3}=2t$
\end{expressions}

\exo[Plusieurs fractions (2)]
Résoudre :
\begin{enumerate}[label=\alph*),itemsep=6pt]
  \item $\dfrac{y}{2}+\dfrac{3y}{4}=\dfrac{y}{3}-22$
  \item $\dfrac{y+7}{3}=\dfrac{y-1}{2}-\dfrac{y}{4}$
  \item $\dfrac{z-4}{3}-\dfrac{z-6}{12}=\dfrac{2z-15}{9}$
  \item $\dfrac{3z+2}{5}-\dfrac{2z-1}{2}=\dfrac{3-6z}{10}$
\end{enumerate}

\partiefiche{Fractions et problèmes}

\exo[Fraction d'une quantité]
Calculer :
\begin{enumerate}[label=\alph*)]
  \item les trois quarts de $120$~€ ;
  \item les deux cinquièmes de $35$~kg ;
  \item les cinq sixièmes de $90$~min ;
  \item les deux tiers des trois quarts des $48$ élèves d'un niveau.
\end{enumerate}

\exo[Le festival]
Un festival peut accueillir $1~200$ spectateurs. Les cinq sixièmes des places sont
vendues en prévente, puis les trois quarts des places restantes sur place. Combien de
places restent invendues ?

\exo[L'anniversaire de Lina]
Lina reçoit $120$~€ pour son anniversaire. Elle en dépense le tiers en jeux vidéo, puis
les trois huitièmes de ce qui lui reste en vêtements.
\begin{enumerate}[label=\alph*)]
  \item Quelle fraction de la somme reçue lui reste-t-il ?
  \item Combien a-t-elle dépensé pour chaque achat ? Combien lui reste-t-il ?
\end{enumerate}

\exo[Le budget d'une association]
Une association consacre les deux cinquièmes de son budget annuel au matériel, puis
les trois quarts de ce qui reste aux sorties. Il reste alors $270$~€ pour les frais de
fonctionnement.
\begin{enumerate}[label=\alph*)]
  \item Quelle fraction du budget représentent les sorties ? les frais de
        fonctionnement ?
  \item Quel est le budget annuel de l'association ?
\end{enumerate}

\partiefiche{Puissances}

\exo[Calculer des puissances]
Calculer :
\begin{expressions}[1.6]
  $A=3^{4}$ & $B=(-2)^{4}$\\
  $C=4^{-2}$ & $D=(-3)^{-3}$\\
  $E=10^{-3}$ & $F=25^{0}$\\
  $G=7^{-1}$ & $H=(-9)^{1}$
\end{expressions}

\exo[Écriture décimale]
Donner l'écriture décimale de :
\begin{expressions}[1.6]
  $A=10^{-3}$ & $B=10^{-7}$\\
  $C=10^{5}$ & $D=10^{0}$
\end{expressions}

\exo[Puissances et priorités]
Calculer :
\begin{expressions}[1.8]
  $A=2+5^{2}$ & $B=(2+5)^{2}$\\
  $C=2-5^{2}$ & $D=3^{3}+2^{4}$\\
  $E=2^{4}\times 3^{2}$ & $F=(-3)^{3}+2^{5}$
\end{expressions}

\exo[Puissances et priorités (suite)]
Calculer :
\begin{expressions}[2.2]
  $A=(3^{3}-4^{2})\times 10^{2}$ & $B=(-3)^{2}+4^{2}\times 10^{-1}$\\
  $C=(-4)^{2}-10^{2}\div 5^{2}$ & $D=10^{2}\times(-1)^{7}-(-2)^{5}$
\end{expressions}

\exo[Des calculs plus longs]
Calculer :
\begin{expressions}[2.6]
  $A=20-3^{2}\times(6-2\times 4)^{2}$ & $B=\dfrac{5^{3}-5\times 4^{2}}{3^{4}-6\times 3^{2}}$\\
  $C=\dfrac{6\times 7-2^{5}}{5\times 2^{3}-(3\times 2)^{2}}$ &
\end{expressions}

\exo[Des résultats en fraction]
Écrire sous forme de fraction simplifiée :
\begin{expressions}[2.6]
  $A=(-3)^{-2}+2\times 10^{-1}$ & $B=\dfrac{3^{2}-(-2)^{3}}{10^{-1}}$\\
  $C=2^{-2}-10^{-1}$ & $D=(2^{5}-4^{2})\times(-4)^{-2}$
\end{expressions}

\partiefiche{Règles de calcul sur les puissances}

\exo[Produits]
Écrire sous la forme $a^{n}$ :
\begin{expressions}[1.8]
  $A=3^{4}\times 3^{7}$ & $B=(-5)^{3}\times(-5)^{5}$\\
  $C=7^{2}\times 7^{-6}$ & $D=x^{3}\times x\times x^{5}$
\end{expressions}

\exo[Quotients]
Écrire sous la forme $a^{n}$ ($x$ et $y$ sont non nuls) :
\begin{expressions}[2.4]
  $A=\dfrac{6^{9}}{6^{4}}$ & $B=\dfrac{(-3)^{8}}{(-3)^{5}}$\\
  $C=\dfrac{y^{12}}{y^{9}}$ & $D=\dfrac{2^{-4}}{2^{3}}$\\
  $E=\dfrac{4^{-1}}{4^{6}}$ & $F=\dfrac{x^{-5}}{x^{-9}}$
\end{expressions}

\exo[Produits et quotients]
Écrire sous la forme $a^{n}$ ($x$ et $y$ sont non nuls) :
\begin{expressions}[2.6]
  $A=\dfrac{3^{10}\times 3^{4}}{3^{5}\times 3^{6}}$ &
  $B=\dfrac{x^{5}}{x^{7}\times x^{4}}$\\
  $C=\dfrac{2^{4}\times 2^{7}}{2^{3}\times 2^{-2}}$ &
  $D=\dfrac{(y^{2})^{5}}{y^{3}\times y^{4}}$
\end{expressions}

\exo[Toutes les règles]
Simplifier, puis calculer :
\[
  A=\dfrac{(5^{2})^{-2}\times 3}{(5\times 3)^{-3}}
  \qquad
  B=\dfrac{3^{4}\times 10^{5}\times 3^{-4}\times 10^{2}}{(3^{3})^{-1}\times(10^{2})^{-4}}
\]

\exo[Calculer et simplifier]
Simplifier ($b$ est non nul) :
\[
  A=\dfrac{2^{6}\times 10^{-3}\times(2^{3})^{2}\times(10^{-1})^{4}}{2^{10}\times 10^{-5}}
  \qquad
  B=\dfrac{b^{2}\times(b^{-3})^{3}}{(b^{4})^{2}}
\]

\partiefiche{Notation scientifique}

\exo[Entiers et décimaux]
Donner la notation scientifique de :
\begin{expressions}[1.6]
  $A=274~000$ & $B=0{,}000~61$\\
  $C=38~200~000$ & $D=0{,}000~000~29$\\
  $E=56~300$ & $F=0{,}000~073~08$
\end{expressions}

\exo[Produits par une puissance de 10]
Donner la notation scientifique de :
\begin{expressions}[1.8]
  $A=532\times 10^{3}$ & $B=64\times 10^{-5}$\\
  $C=481{,}7\times 10^{8}$ & $D=0{,}036\times 10^{-4}$
\end{expressions}

\exo[Un calcul jusqu'au bout]
Donner la notation scientifique de
\[ E=\dfrac{35\times 10^{-2}\times 6\times(10^{3})^{3}}{14\times 10^{5}}. \]

\exo[Deux expressions]
On pose
\[
  A=\dfrac{3\times 10^{1500}}{12\times 10^{1502}}
  \quad\text{et}\quad
  B=\dfrac{6{,}4\times 10^{2}-340\times 10^{-1}}{1{,}01\times 10^{2}}.
\]
\begin{enumerate}[label=\arabic*)]
  \item Donner la notation scientifique de $A$.
  \item Montrer que $B$ est un nombre entier.
\end{enumerate}

\exo[La lumière de la Lune]
La lumière parcourt environ $3\times 10^{8}$~m par seconde. La Lune est à environ
$384~000$~km de la Terre.
\begin{enumerate}[label=\arabic*)]
  \item Écrire la distance Terre-Lune en mètres, en notation scientifique.
  \item Combien de temps la lumière renvoyée par la Lune met-elle pour nous parvenir ?
\end{enumerate}

\end{multicols}

\end{document}
